\documentclass[11pt]{article}
\usepackage[margin=1in]{geometry}
\usepackage{amsmath,amssymb,amsthm}
\usepackage{graphicx}
\usepackage{booktabs}
\usepackage{hyperref}
\usepackage{caption}
\usepackage{subcaption}
\usepackage{algorithm}
\usepackage{algpseudocode}
\usepackage{xcolor}
\usepackage{microtype}
\usepackage{float}
\usepackage{parskip}
\usepackage{enumitem}

\newtheorem{theorem}{Theorem}



\hypersetup{colorlinks=true, linkcolor=blue!60!black, citecolor=blue!60!black,
            urlcolor=blue!60!black}

\title{\textbf{Cardinality-Constrained Portfolio Optimization:\\[4pt]
Empirical Computational Hardness on S\&P~500 Data}}
\author{Aashrit Cunchala \and Aymen Marzak \and Adhish Rajan \and Jean SainteBeuve \and Clarys Vignal}
\date{Spring 2026 \quad --- \quad Discrete Probability and Stochastic Processes}

\begin{document}
\maketitle

% ─────────────────────────────────────────────────────────────────────────────
\begin{abstract}
% ─────────────────────────────────────────────────────────────────────────────
We study the problem of selecting a portfolio of exactly $k$ assets from a
universe of $N=462$ S\&P~500 stocks to maximise expected return subject
to a variance constraint, with short selling allowed.  Without a cardinality
restriction the problem reduces to the classical Markowitz mean-variance
optimisation, which admits a closed-form solution via the two-fund separation
theorem.  Adding the constraint $|\operatorname{supp}(\mathbf{w})|=k$ makes the
search space $\binom{N}{k}$, which grows super-exponentially in $N$ and renders
exact enumeration infeasible for realistic values.  We exploit the closed-form
inner solution to evaluate each candidate $k$-subset exactly and efficiently, and
we compare four approaches: the Markowitz upper bound (no cardinality), brute-force
enumeration, a greedy forward-selection heuristic, and simulated annealing
(SA)---interpreted as a Markov chain on the space of $k$-subsets.  We further
study the standard convex relaxation of the cardinality constraint, which replaces
the L0 penalty $\lambda\|\mathbf{w}\|_0$ with an L1 penalty $\lambda\|\mathbf{w}\|_1$
in the objective (LASSO-type regularisation) and is solvable as a second-order cone
programme.  Our experiments quantify (i)~the exponential blow-up in search space
and runtime, (ii)~the quality gap between heuristics and the exact optimum,
(iii)~SA convergence behaviour as a function of chain length, $k$, and $n$,
(iv)~the cost of the cardinality constraint as $k$ grows toward~$N$,
(v)~the sensitivity of algorithmic hardness to the tightness of the risk constraint,
(vi)~how much of the search space must be sampled to find a near-optimal portfolio,
and (vii)~the regularisation path and long-only floor of the L1 relaxation.
\end{abstract}

% ─────────────────────────────────────────────────────────────────────────────
\section{Problem Formulation}
\label{sec:formulation}
% ─────────────────────────────────────────────────────────────────────────────

Let $\boldsymbol{\mu}\in\mathbb{R}^N$ be the vector of annualised expected
returns, $\boldsymbol{\Sigma}\in\mathbb{R}^{N\times N}$ the annualised
covariance matrix (positive definite), $\sigma_0^2>0$ a risk budget, and
$k\leq N$ a prescribed portfolio size.  Portfolio weights
$\mathbf{w}\in\mathbb{R}^N$ satisfy the budget constraint $\mathbf{1}^\top
\mathbf{w}=1$; \textbf{short selling is allowed}, i.e.\ $w_i<0$ is permitted.

\subsection{Continuous Relaxation: Markowitz Mean-Variance (Polynomial Time)}

Dropping the cardinality constraint gives the classical Markowitz problem
\cite{Markowitz1952}:
\begin{equation}
  R^*_{\text{MV}}
  \;=\;
  \max_{\mathbf{w}\in\mathbb{R}^N}\; \boldsymbol{\mu}^\top\mathbf{w}
  \quad\text{s.t.}\quad
  \mathbf{w}^\top\boldsymbol{\Sigma}\mathbf{w}\leq\sigma_0^2,\;
  \mathbf{1}^\top\mathbf{w}=1.
  \label{eq:mv}
\end{equation}
Because $\boldsymbol{\Sigma}\succ 0$ and the feasible set is a convex ellipsoid,
problem~\eqref{eq:mv} is strictly convex and admits a \emph{closed-form} solution.

\paragraph{Closed-form derivation.}
Define the scalars
\begin{equation}
  A \;=\; \mathbf{1}^\top\boldsymbol{\Sigma}^{-1}\mathbf{1}, \quad
  B \;=\; \mathbf{1}^\top\boldsymbol{\Sigma}^{-1}\boldsymbol{\mu}, \quad
  C \;=\; \boldsymbol{\mu}^\top\boldsymbol{\Sigma}^{-1}\boldsymbol{\mu}, \quad
  D \;=\; AC - B^2 > 0.
  \label{eq:abcd}
\end{equation}
Because higher expected return requires a larger variance, the variance constraint
is \emph{active} at the optimum.  The Lagrangian with multipliers $\lambda>0$
(variance) and $\nu$ (budget) is
\begin{equation}
  \mathcal{L}(\mathbf{w},\lambda,\nu)
  \;=\;
  \boldsymbol{\mu}^\top\mathbf{w}
  \;-\; \lambda\bigl(\mathbf{w}^\top\boldsymbol{\Sigma}\mathbf{w} - \sigma_0^2\bigr)
  \;-\; \nu\bigl(\mathbf{1}^\top\mathbf{w} - 1\bigr).
  \label{eq:lagrangian}
\end{equation}
The first-order condition $\partial\mathcal{L}/\partial\mathbf{w}=\mathbf{0}$ gives
\begin{equation}
  \boldsymbol{\mu} - 2\lambda\boldsymbol{\Sigma}\mathbf{w} - \nu\mathbf{1}
  \;=\; \mathbf{0}
  \quad\Longrightarrow\quad
  \mathbf{w} \;=\; \frac{1}{2\lambda}\,\boldsymbol{\Sigma}^{-1}
    \bigl(\boldsymbol{\mu}-\nu\mathbf{1}\bigr).
  \label{eq:foc}
\end{equation}
Equation~\eqref{eq:foc} is the \emph{two-fund separation theorem}: every efficient
portfolio is a linear combination of $\boldsymbol{\Sigma}^{-1}\boldsymbol{\mu}$
(tangency direction) and $\boldsymbol{\Sigma}^{-1}\mathbf{1}$ (minimum-variance
direction).  Writing $\mathbf{w}^*=\alpha\,\boldsymbol{\Sigma}^{-1}\boldsymbol{\mu}
+\beta\,\boldsymbol{\Sigma}^{-1}\mathbf{1}$ and imposing the budget and target-return
constraints at level $\mu^*$:
\begin{align}
  \mathbf{1}^\top\mathbf{w}^*=1 &\;\Longrightarrow\; \alpha B + \beta A \;=\; 1,
  \label{eq:budget_cstr}\\
  \boldsymbol{\mu}^\top\mathbf{w}^*=\mu^*
  &\;\Longrightarrow\; \alpha C + \beta B \;=\; \mu^*.
  \label{eq:return_cstr}
\end{align}
Solving the $2\times2$ system~\eqref{eq:budget_cstr}--\eqref{eq:return_cstr}
(determinant $D=AC-B^2>0$):
\begin{equation}
  \alpha \;=\; \frac{A\mu^*-B}{D}, \qquad \beta \;=\; \frac{C-B\mu^*}{D}.
  \label{eq:alphabeta}
\end{equation}
Substituting~\eqref{eq:alphabeta} into $\mathbf{w}^*=\alpha\boldsymbol{\Sigma}^{-1}
\boldsymbol{\mu}+\beta\boldsymbol{\Sigma}^{-1}\mathbf{1}$ yields the optimal weights:
\begin{equation}
  \mathbf{w}^* \;=\; \frac{1}{D}\Bigl[
    \bigl(C - B\mu^*\bigr)\,\boldsymbol{\Sigma}^{-1}\mathbf{1}
    \;+\;
    \bigl(A\mu^* - B\bigr)\,\boldsymbol{\Sigma}^{-1}\boldsymbol{\mu}
  \Bigr].
  \label{eq:mvweights}
\end{equation}
The frontier variance is
\begin{equation}
  \mathbf{w}^{*\top}\boldsymbol{\Sigma}\mathbf{w}^*
  \;=\; \alpha^2 C + 2\alpha\beta B + \beta^2 A
  \;=\; \frac{1}{A} + \frac{A}{D}\!\left(\mu^* - \frac{B}{A}\right)^{\!2},
  \label{eq:frontier}
\end{equation}
which is the \emph{efficient frontier} parabola in $(\mu^*,\sigma^2)$ space,
with global minimum-variance portfolio at $\sigma^2_{\min}=1/A$, attained when
$\mu^*=B/A$.  Setting~\eqref{eq:frontier} equal to $\sigma_0^2$ and solving
for $\mu^*$ (taking the upper branch, since larger $\mu^*$ requires higher variance):
\begin{equation}
  R^*_{\text{MV}} \;=\; \frac{B}{A} + \sqrt{\frac{D}{A}\left(\sigma_0^2
  - \frac{1}{A}\right)}.
  \label{eq:mvreturn}
\end{equation}
Problem~\eqref{eq:mv} is infeasible if and only if $\sigma_0^2 < 1/A$.

\subsection{Cardinality-Constrained Problem (NP-Hard)}

Fix a cardinality $k \leq N$.  For each subset $S \subseteq [N]$ with
$|S|=k$, define the \emph{inner Markowitz problem} restricted to $S$:
\begin{equation}
  v(S) \;=\;
  \max_{\mathbf{w}_S\in\mathbb{R}^k}\;
    \boldsymbol{\mu}_S^\top\mathbf{w}_S
  \quad\text{s.t.}\quad
  \mathbf{w}_S^\top\boldsymbol{\Sigma}_S\mathbf{w}_S\leq\sigma_0^2,\;
  \mathbf{1}^\top\mathbf{w}_S=1,
  \label{eq:inner}
\end{equation}
where $\boldsymbol{\mu}_S$ and $\boldsymbol{\Sigma}_S$ are the restrictions of
$\boldsymbol{\mu}$ and $\boldsymbol{\Sigma}$ to indices in $S$.  By
equations~\eqref{eq:mvreturn}--\eqref{eq:mvweights} applied to the $k$-dimensional
sub-problem, $v(S)$ is available in \emph{closed form} for any $S$.

The cardinality-constrained problem is:
\begin{equation}
  \max_{\substack{S\subseteq[N] \\ |S|=k}}\; v(S).
  \label{eq:card}
\end{equation}
Problem~\eqref{eq:card} is NP-hard.  This can be established via a polynomial
reduction from the Maximum Weight Independent Set problem: given a weighted graph
$G=(V,E)$ with $|V|=N$, one constructs a mean-variance instance on $N$ assets
whose optimal $k$-subset of maximum expected return corresponds to a maximum-weight
independent set of $G$; the reduction runs in polynomial time
\cite{GareyJohnson1979}.  Consequently, no polynomial-time exact algorithm for
problem~\eqref{eq:card} is expected to exist (unless $\mathsf{P}=\mathsf{NP}$),
and the exponential growth of $\binom{N}{k}$ confirmed empirically in
Experiment~A (\S\ref{sec:expA}) is a direct manifestation of this hardness.



\subsection{L1 Relaxation (LASSO-Type Surrogate)}
\label{sec:lasso_theory}

Because problem~\eqref{eq:card} is NP-hard, a natural approach is to replace
the intractable $\ell_0$ constraint with a convex surrogate.  The standard choice
is to move the $\ell_0$ penalty into the objective and replace it with the $\ell_1$
norm \cite{Tibshirani1996}:
\begin{equation}
  \max_{\mathbf{w}\in\mathbb{R}^n}\;
    \boldsymbol{\mu}^\top\mathbf{w} - \lambda\|\mathbf{w}\|_1
  \quad\text{s.t.}\quad
  \mathbf{w}^\top\boldsymbol{\Sigma}\mathbf{w}\leq\sigma_0^2,\;
  \mathbf{1}^\top\mathbf{w}=1,
  \label{eq:lasso}
\end{equation}
where $\lambda\geq 0$ controls the trade-off between expected return and gross
exposure $\|\mathbf{w}\|_1 = \sum_i|w_i|$.  Problem~\eqref{eq:lasso} is a
second-order cone programme (SOCP) and is solvable in polynomial time.

\paragraph{KKT conditions and soft-thresholding.}
The Lagrangian of~\eqref{eq:lasso} includes a non-smooth term from $\|\mathbf{w}\|_1$.
The first-order KKT stationarity condition at the optimum $\mathbf{w}^*$ is
\begin{equation}
  \mu_i \;-\; 2\gamma\bigl(\boldsymbol{\Sigma}\mathbf{w}^*\bigr)_i \;-\; \nu
  \;\in\; \lambda\,\partial|w_i^*|,
  \label{eq:lasso_kkt}
\end{equation}
where $\gamma\geq 0$ is the Lagrange multiplier on the variance constraint and $\nu$
enforces the budget.  The subdifferential $\partial|w_i^*|$ equals $\{\operatorname{sign}(w_i^*)\}$ for
$w_i^*\neq 0$ and $[-1,+1]$ for $w_i^*=0$.  Condition~\eqref{eq:lasso_kkt} yields
a soft-thresholding rule: stock~$i$ is \emph{excluded} ($w_i^*=0$) if and only if
\begin{equation}
  \bigl|\mu_i - 2\gamma\bigl(\boldsymbol{\Sigma}\mathbf{w}^*\bigr)_i - \nu\bigr|
  \;\leq\; \lambda,
  \label{eq:exclusion}
\end{equation}
i.e.\ its risk-adjusted excess return lies within the penalty band $[-\lambda,\lambda]$.
This is the portfolio analogue of the LASSO soft-thresholding rule in regression.
Unlike regression LASSO, however, the threshold in~\eqref{eq:exclusion} is
asset-specific through the covariance term $(\boldsymbol{\Sigma}\mathbf{w}^*)_i$,
so there is no closed-form solution and a numerical solver is required.

\paragraph{The long-only floor.}
A structural limitation of the $\ell_1$ penalty arises from the budget constraint
$\mathbf{1}^\top\mathbf{w}=1$.  For a portfolio with only long positions
($w_i\geq 0$ for all $i$), the gross exposure equals the budget:
$\|\mathbf{w}\|_1=\mathbf{1}^\top\mathbf{w}=1$.  Once~\eqref{eq:lasso} has
eliminated all short positions, increasing $\lambda$ further leaves $\|\mathbf{w}\|_1$
fixed at its minimum of $1$ for any allocation of the budget.  The $\ell_1$
penalty therefore cannot distinguish among long-only portfolios with different
cardinalities: all achieve the same penalty $\lambda\cdot 1$.  Consequently,
the LASSO path reaches a \emph{long-only floor}: a minimum cardinality
$k_{\min}^{\ell_1}$ beyond which further sparsification is impossible via $\ell_1$
regularisation alone.  The L1 penalty confounds \emph{gross-exposure reduction}
(eliminating short positions) with \emph{cardinality reduction} (selecting fewer
stocks), conflating two structurally different operations into a single penalty.
This limitation is studied empirically in Experiment~G (\S\ref{sec:expG}).

% ─────────────────────────────────────────────────────────────────────────────
\section{Data}
\label{sec:data}
% ─────────────────────────────────────────────────────────────────────────────

We use daily adjusted-close returns for all $N=462$ S\&P~500 constituents
available for the full sample period.  Tickers with more than 20\% missing
observations are excluded; the remaining rows are kept only if complete across
all retained tickers, yielding a balanced panel of $T=2{,}854$ trading days
(January 2015 to May 2026) with no missing values.

From the return matrix $\mathbf{R}\in\mathbb{R}^{T\times N}$:
\[
  \hat{\boldsymbol{\mu}} = 252\cdot\bar{\mathbf{r}},
  \qquad
  \hat{\boldsymbol{\Sigma}} = 252\cdot\widehat{\mathrm{Cov}}(\mathbf{R}).
\]

\subsection{Descriptive Statistics}
\label{sec:descriptive}

Figure~\ref{fig:data_overview} provides an overview of the return data.
The top panel plots the cumulative wealth of six representative stocks
(one per sector) over the full sample, illustrating the wide dispersion
in long-run performance and the heterogeneity of risk.
The middle panel displays the cross-sectional risk--return tradeoff
for all 462 stocks together with the equal-weight portfolio.
The bottom panel shows the histogram of daily returns for AAPL against the
fitted Normal distribution, illustrating the fat tails characteristic of
financial return series.

\begin{figure}[H]
  \centering
  \includegraphics[width=\linewidth]{fig0_data_overview.png}
  \caption{Data overview.
    \textit{Top}: cumulative wealth (initial \$1) for six representative stocks.
    \textit{Middle}: annualised risk--return scatter for all 462 S\&P~500 stocks;
      the red star marks the equal-weight portfolio.
    \textit{Bottom}: daily return histogram for AAPL vs.\ a fitted Normal
      distribution, highlighting negative skewness and excess kurtosis.}
  \label{fig:data_overview}
\end{figure}

Table~\ref{tab:descriptive} summarises the cross-sectional distribution of
key statistics for all 462 stocks.  Annualised returns range from $-1.56\%$
(VTRS) to $65.67\%$ (NVDA) with a mean of $16.19\%$; annualised volatilities
from $17.76\%$ to $68.45\%$ with a mean of $31.29\%$.  Daily return
distributions exhibit strong negative skewness and pronounced fat tails:
the median excess kurtosis is $11.43$, far exceeding the $0$ predicted by
the Normal model.  These features motivate the use of in-sample variance as
the risk measure, as the tail properties are too heterogeneous across stocks
for a simple parametric treatment.

\begin{table}[ht]
\centering
\caption{Cross-sectional descriptive statistics for daily returns of the
  $N=462$ S\&P~500 stocks (annualised where indicated).}
\label{tab:descriptive}
\begin{tabular}{lrrrrrr}
\toprule
Statistic & Min & Q25 & Median & Mean & Q75 & Max \\
\midrule
Ann.\ return (\%)     & $-1.56$ & $10.86$ & $14.95$ & $16.19$ & $19.59$ & $65.67$ \\
Ann.\ volatility (\%) & $17.76$ & $25.64$ & $29.42$ & $31.29$ & $35.34$ & $68.45$ \\
Daily skewness        & $-6.76$ & $-0.39$ & $-0.06$ & $-0.07$ & $0.19$  & $11.27$ \\
Excess kurtosis       & $3.28$  & $8.26$  & $11.43$ & $15.53$ & $16.07$ & $357.27$ \\
\bottomrule
\end{tabular}
\end{table}

\subsection{Sample Covariance Matrix: Invertibility and Spectral Properties}
\label{sec:rmt}

\paragraph{Invertibility and condition number.}
The sample covariance matrix $\hat{\boldsymbol{\Sigma}}$ is positive definite if
and only if the return matrix $\mathbf{R}$ has full column rank, which requires
$T > N$.  With $T=2{,}854$ and $N=462$ we have $T/N\approx 6$, so
$\hat{\boldsymbol{\Sigma}}$ is invertible but can be poorly conditioned.  The
condition number $\kappa(\hat{\boldsymbol{\Sigma}})=\lambda_{\max}/\lambda_{\min}$
measures this: a large $\kappa$ means the inverse amplifies the directions
corresponding to the smallest eigenvalues, assigning extreme weights to those
eigendirections.  In high-dimensional financial datasets, the smallest eigenvalues
are typically dominated by estimation noise, making the resulting Markowitz
weights unstable and the expected-return estimate unreliable.

\paragraph{Marchenko--Pastur law.}
Under the null hypothesis that returns are i.i.d.\ Gaussian with variance
$\sigma^2$, the empirical spectral distribution of $\hat{\boldsymbol{\Sigma}}$
converges (as $T,N\to\infty$ with $T/N\to q>1$) to the
\emph{Marchenko--Pastur law} \cite{MarchenkoPastur1967}:
\begin{equation}
  f_{\mathrm{MP}}(\lambda)
  \;=\;
  \frac{q}{2\pi\sigma^2}
  \cdot
  \frac{\sqrt{(\lambda_+ - \lambda)(\lambda - \lambda_-)}}{\lambda},
  \quad \lambda\in[\lambda_-,\,\lambda_+],
  \label{eq:mp}
\end{equation}
with bulk edges
\begin{equation}
  \lambda_{\pm} \;=\; \sigma^2\!\left(1 \pm \frac{1}{\sqrt{q}}\right)^{\!2}.
  \label{eq:mpedges}
\end{equation}
Eigenvalues within $[\lambda_-,\lambda_+]$ are statistically indistinguishable
from pure noise; only eigenvalues \emph{above} $\lambda_+$ represent genuine
cross-asset correlations.

For our dataset with $q=T/N\approx 6$, equation~\eqref{eq:mpedges} gives
$\lambda_-/\sigma^2\approx 0.36$ and $\lambda_+/\sigma^2\approx 1.96$.
A large fraction of the 462 eigenvalues of $\hat{\boldsymbol{\Sigma}}$ therefore
fall within the Marchenko--Pastur bulk and carry no reliable information about
the true covariance structure.  Inverting $\hat{\boldsymbol{\Sigma}}$ assigns
large weights ($\sim 1/\lambda_{\min}$) to these noise-contaminated eigenvectors,
which explains the inflated Markowitz upper bound of $R^*_{\text{MV}}=120.24\%$
reported in Section~\ref{sec:algorithms} and Experiment~D: this figure reflects
in-sample optimisation under a noisy inverse, not a reliable out-of-sample forecast.

\paragraph{Implication for cardinality-constrained experiments.}
The inner problems use $k\times k$ sub-matrices.  For all tractable experiments
($k\leq 30$) the ratio $T/k>90$, placing those sub-problems in a well-conditioned
regime where the Marchenko--Pastur bulk is narrow and eigenvalue estimation is
reliable.  The quality comparisons in Experiment~B are therefore unaffected by
the full-$N$ conditioning issue.

\paragraph{Risk budget.}
We set
\begin{equation}
  \sigma_0^2 \;=\; \mathrm{Quantile}_{25\%}\!\left(\operatorname{diag}
    \hat{\boldsymbol{\Sigma}}\right),
  \label{eq:riskbudget}
\end{equation}
i.e.\ the 25th percentile of individual stock variances, giving
$\sigma_0 = 25.64\%$.  This ensures the inner problem~\eqref{eq:inner} is
feasible for a substantial fraction of $k$-subsets: feasibility requires
$\sigma_0^2 \geq 1/A_S$, which holds whenever the subset contains at least one
sufficiently low-variance stock and weights can be chosen freely (short selling
helps by allowing variance-reducing hedges).

% ─────────────────────────────────────────────────────────────────────────────
\section{Algorithms}
\label{sec:algorithms}
% ─────────────────────────────────────────────────────────────────────────────

\subsection{Markowitz Upper Bound}

We solve~\eqref{eq:mv} via the closed-form~\eqref{eq:mvreturn}--\eqref{eq:mvweights}
with all $N=462$ assets.  The result $R^*_{\text{MV}}$ upper-bounds every
cardinality-constrained solution.  Cost: $O(N^3)$ for the matrix inversion
$\boldsymbol{\Sigma}^{-1}$.

\subsection{Brute-Force Enumeration}

For small $n$ and $k$ we enumerate all $\binom{n}{k}$ subsets and evaluate
$v(S)$ via the closed-form inner solution for each.  Since each evaluation is
$O(k^3)$, the total cost is $O\!\left(k^3\binom{n}{k}\right)$; it is
feasible only for $\binom{n}{k}\lesssim 10^5$.

\subsection{Greedy Forward Selection}

Starting from $S=\emptyset$, we greedily add the asset that maximises
$v(S\cup\{i\})$ among all remaining candidates, repeating $k$ times.
Cost: $O(k\cdot n)$ closed-form evaluations.

\begin{algorithm}[h]
\caption{Sequential Greedy}
\begin{algorithmic}[1]
\State $S \leftarrow \emptyset$,\quad $C \leftarrow \{1,\ldots,N\}$
\For{$t = 1,\ldots,k$}
  \State $i^* \leftarrow \arg\max_{i\in C}\; v(S\cup\{i\})$
  \If{no feasible $i^*$ exists} \textbf{break} \EndIf
  \State $S \leftarrow S\cup\{i^*\}$,\quad $C \leftarrow C\setminus\{i^*\}$
\EndFor
\State \Return $S$, $v(S)$
\end{algorithmic}
\end{algorithm}

\subsection{Simulated Annealing as a Markov Chain}
\label{sec:sa}

We model the search as a \emph{Markov chain} on the discrete state space
$\mathcal{X} = \binom{[N]}{k}$ of all $k$-subsets, with
$|\mathcal{X}|=\binom{N}{k}$.  The transition kernel at state $S$ proposes a
\textbf{swap}: draw $i\sim\mathrm{Uniform}(S)$ and $j\sim\mathrm{Uniform}
([N]\setminus S)$, forming the candidate $S'=(S\setminus\{i\})\cup\{j\}$.
The move is accepted with the Metropolis--Hastings probability
\[
  \alpha(S,S') \;=\; \min\!\left(1,\,e^{\bigl(v(S')-v(S)\bigr)/T_s}\right),
\]
where the temperature decreases geometrically:
\begin{equation}
  T_s = T_0\!\left(\frac{T_{\min}}{T_0}\right)^{s/M},
  \quad s=0,1,\ldots,M,
  \label{eq:cooling}
\end{equation}
with $M=2{,}000$, $T_0=0.004$, $T_{\min}=5\times10^{-5}$.

The proposal kernel is symmetric (swap is reversible), so the chain satisfies
detailed balance with respect to the Gibbs measure $\pi_T(S)\propto e^{v(S)/T}$
at each fixed temperature.  As $T\to 0$ the chain concentrates on the global
maximisers of $v$.  We run $R=3$ independent restarts and keep the global best.

\paragraph{Computational cost per step.}
Since short selling is permitted and the only constraint besides the variance
bound is the budget equation $\mathbf{1}^\top\mathbf{w}_S=1$, the inner problem
for each candidate subset reduces to the classical Markowitz problem and admits
the closed-form solution~\eqref{eq:mvreturn}--\eqref{eq:mvweights}.  Each
evaluation of $v(S)$ therefore requires a single matrix inversion ($O(k^3)$) and
a handful of dot products; there is no iterative solver overhead.  For
$k\leq 20$ this amounts to at most $\sim\!8{,}000$ floating-point operations,
so the computational bottleneck is entirely the combinatorial outer search.

\subsection{Mixing-Time Analysis of the Swap Chain}
\label{sec:mixing}

The observed fast convergence of SA ($O(10^3)$ steps even when
$|\mathcal{X}|=\binom{N}{k}\approx 10^{20}$ for $k=10$) follows from
a quantitative spectral-gap bound on the underlying Markov chain.

\paragraph{State space and proposal kernel.}
At any fixed temperature $T_s$, the SA chain is a reversible Markov
chain on $\mathcal{X}=\binom{[N]}{k}$ with stationary distribution
$\pi_{T_s}(S)\propto e^{v(S)/T_s}$.
The proposal kernel draws a uniformly random \emph{swap}
$(i,j)\in S\times([N]\setminus S)$, so each state $S$ has exactly
\begin{equation}
  \deg(S) = k(N-k)
  \label{eq:degree}
\end{equation}
neighbours, and the underlying undirected graph is the
\emph{Johnson graph} $J(N,k)$ \cite{Godsil1993}.

\paragraph{Spectral gap of the uniform swap chain.}
Consider first the \emph{uniform} chain obtained by setting $T_s\to\infty$
(all swaps accepted); its stationary distribution is uniform on
$\mathcal{X}$.
The eigenvalues of the transition matrix of $J(N,k)$ are known from
the theory of association schemes \cite{Godsil1993}: for $r=0,1,\ldots,k$,
\begin{equation}
  \lambda_r
  \;=\;
  1 - \frac{r(N+1-r)}{k(N-k)}.
  \label{eq:jkeval}
\end{equation}
The spectral gap $\gamma^* = 1-\lambda_1$ is therefore
\begin{equation}
  \boxed{
  \gamma^* \;=\; \frac{N}{k(N-k)}.
  }
  \label{eq:gap}
\end{equation}
For $N=462$, $k=10$: $\gamma^*=462/(10\cdot 452)\approx 0.102$;
for $k=50$: $\gamma^*\approx 0.024$.

\paragraph{Mixing-time bound.}
For a reversible chain with spectral gap $\gamma^*$ and
$|\mathcal{X}|=\binom{N}{k}$, the standard mixing-time theorem
\cite{LevinPeres2017} gives
\begin{equation}
  \tau_{\mathrm{mix}}(\delta)
  \;\leq\;
  \frac{1}{\gamma^*}\!\left[\log|\mathcal{X}|+\log\frac{1}{\delta}\right]
  \;=\;
  \frac{k(N-k)}{N}
  \left[\log\binom{N}{k}+\log\frac{1}{\delta}\right].
  \label{eq:mixbound}
\end{equation}
Using $\log\binom{N}{k}\leq k\log(eN/k)$:
\begin{equation}
  \tau_{\mathrm{mix}}(\delta)
  \;=\;
  O\!\left(\frac{k^2(N-k)}{N}\log\frac{N}{k}\right).
  \label{eq:mixO}
\end{equation}
Plugging in $N=462$, $k=10$, $\delta=0.05$:
\[
  \tau_{\mathrm{mix}} \;\lesssim\;
  \frac{10^2\cdot 452}{462}\cdot\log\!\left(\frac{462}{10}\right)
  \;\approx\; 748\text{ steps},
\]
consistent with the empirical plateau at $\sim\!500$ steps
(Experiment~C, \S\ref{sec:expC}).
For $k=50$: the bound gives $\approx 9{,}900$ steps, explaining the
slower convergence observed for large cardinalities.

\paragraph{Effect of the Metropolis acceptance step.}
At finite temperature $T_s$, a proposed move is accepted with
probability $\alpha(S,S')=\min(1,e^{(v(S')-v(S))/T_s})\geq e^{-\Delta/T_s}$,
where $\Delta=\max_{S,S'\text{ adjacent}}|v(S)-v(S')|$ is the maximum
energy change over a single swap.
The spectral gap of the Metropolis chain satisfies \cite{LawlerSokal1988}
\begin{equation}
  \gamma^*_{\mathrm{MH}}
  \;\geq\;
  e^{-\Delta/T_s}\,\gamma^*_{\mathrm{unif}},
  \label{eq:mhgap}
\end{equation}
so at high temperatures (early in the SA schedule, large $T_s$),
the gap is close to $\gamma^*=N/[k(N-k)]$ and mixing is rapid.
As $T_s\to 0$ the gap contracts, but the chain has already
concentrated near the optimal region of $\mathcal{X}$, so further
refinement requires only local moves.

\paragraph{Implication for the cooling schedule.}
Equation~\eqref{eq:mixO} establishes that $O(k^2\log(N/k))$ steps
suffice for the \emph{uniform} chain to mix, independently of
$|\mathcal{X}|$.
The super-exponential cardinality of $\mathcal{X}$ does not appear in
this bound; what matters is the degree $k(N-k)$, which keeps the spectral
gap~\eqref{eq:gap} bounded away from zero.
For the SA schedule~\eqref{eq:cooling} with $M=2{,}000$ steps,
the total chain length comfortably exceeds $\tau_{\mathrm{mix}}$ in
the high-temperature phase, giving SA its empirically fast convergence
to near-optimal portfolios even for $N=462$.


% ─────────────────────────────────────────────────────────────────────────────
\section{Experiments}
\label{sec:experiments}
% ─────────────────────────────────────────────────────────────────────────────

All experiments use the risk budget~\eqref{eq:riskbudget} ($\sigma_0=25.64\%$).

\paragraph{Experimental design: tractability and fair comparison.}
Brute-force enumeration is the only method that produces a provably exact optimum,
but its cost $O\!\left(k^3\binom{n}{k}\right)$ limits it to instances where
$\binom{n}{k}\lesssim 10^5$, i.e.\ $n\lesssim 30$ for $k=5$ (see
Experiment~A).  Any experiment that uses brute force as a benchmark must
therefore restrict \emph{all} competing methods --- including the Markowitz
upper bound --- to the \emph{same} $n$-stock sub-universe; otherwise the
comparison conflates the effect of the algorithm with the effect of working on
a different instance.  For experiments that do not require an exact reference
(Experiments C--D), the full $N=462$ universe is used for both greedy and SA.

\subsection{Experiment A: Computational Scaling}
\label{sec:expA}

Figure~\ref{fig:scaling} (left) plots $\binom{n}{k}$ as a function of $n$ for
$k\in\{5,10,20\}$ on a log scale, illustrating the super-exponential growth.
Figure~\ref{fig:scaling} (right) records actual brute-force runtimes for $k=5$.
Measured times are: $n=10$: $252$ subsets, ${<}0.1$\,s; $n=15$: $3{,}003$
subsets, $0.1$\,s; $n=20$: $15{,}504$ subsets, $0.5$\,s; $n=25$: $53{,}130$
subsets, $1.8$\,s; $n=30$: $142{,}506$ subsets, $4.9$\,s.
A log-log extrapolation predicts infeasibility beyond $n\approx 40$--45 for $k=5$.

\begin{figure}[ht]
  \centering
  \includegraphics[width=\linewidth]{fig2_scaling.png}
  \caption{Search space size $\binom{n}{k}$ (left) and measured brute-force
    runtime (right) as functions of universe size~$n$ for $k=5$.}
  \label{fig:scaling}
\end{figure}

\subsection{Experiment B: Algorithm Quality Comparison}
\label{sec:expB}

For $n=30$, $k=5$ ($\binom{30}{5}=142{,}506$ subsets, within the brute-force
budget), we compare all four methods on the \emph{same} curated sub-universe
of 30 large-cap S\&P~500 stocks spanning six sectors (Technology, Financials,
Healthcare, Energy, Consumer Staples, and Consumer Discretionary).  Crucially,
the Markowitz upper bound is computed on this 30-stock sub-universe --- not the
full $N=462$ universe --- so that all four numbers reflect identical input data
and the comparison is exact.  Table~\ref{tab:quality} summarises the results;
Figure~\ref{fig:quality} shows the bar chart.

\begin{table}[ht]
\centering
\caption{Algorithm quality comparison ($n=30$, $k=5$, $\sigma_0=25.64\%$).
  Gap is relative to brute-force optimal.
  The MV--BF gap measures the price of the cardinality constraint.}
\label{tab:quality}
\begin{tabular}{lcc}
\toprule
Method & Ann.\ return & Gap vs.\ BF optimal \\
\midrule
Markowitz upper bound (no cardinality) & $47.79\%$ & $+12.56\%$\textsuperscript{a} \\
Brute-force (optimal $k=5$)            & $41.79\%$ & $0\%$ \\
Greedy                                 & $40.49\%$ & $3.11\%$ \\
Simulated Annealing                    & $41.79\%$ & $0.00\%$ \\
\bottomrule
\multicolumn{3}{l}{\small\textsuperscript{a}Cardinality gap: cost of restricting to 5 assets.}
\end{tabular}
\end{table}

\begin{figure}[ht]
  \centering
  \includegraphics[width=0.62\linewidth]{fig3_quality_comparison.png}
  \caption{Annualised return achieved by each method ($n=30$, $k=5$).
    Red = Markowitz upper bound; green = brute-force optimum.}
  \label{fig:quality}
\end{figure}

SA matches the brute-force optimum exactly ($41.79\%$, $0.00\%$ gap), while
greedy obtains $40.49\%$ ($3.11\%$ gap).
The $12.56\%$ gap between the Markowitz upper bound ($47.79\%$) and the
brute-force optimal ($41.79\%$) quantifies the \emph{price of the cardinality
constraint} at $k=5$.

\subsection{Experiment C: SA Convergence as a Markov Chain}
\label{sec:expC}

No exact benchmark is needed here; greedy and SA are run on the full $N=462$
universe.  Figure~\ref{fig:sa_convergence} traces the best-so-far return against
Markov chain step, averaged over $R=3$ restarts, for varying $k$ with $n=462$
fixed (left) and varying $n$ with $k=10$ fixed (right).

\begin{figure}[ht]
  \centering
  \includegraphics[width=\linewidth]{fig4_sa_convergence.png}
  \caption{SA convergence curve.  Left: $n=462$ fixed, $k\in\{5,10,20,50\}$
    varied.  Right: $k=10$ fixed, $n\in\{20,50,462\}$ varied.
    Each curve is averaged over 3 restarts.}
  \label{fig:sa_convergence}
\end{figure}

Final SA returns for $n=462$ were: $k=5$: $49.07\%$; $k=10$: $57.72\%$;
$k=20$: $68.24\%$; $k=50$: $84.91\%$.
For $k=10$ varying $n$: $n=20$: $32.68\%$; $n=50$: $46.31\%$; $n=462$: $57.33\%$.

\begin{itemize}[leftmargin=1.5em]
  \item Convergence is rapid in the first $\sim\!500$ steps and plateaus
    thereafter, independent of $n$ or $k$.
  \item Larger $k$ attains higher absolute returns; the convergence
    \emph{rate} is similar across $k$.
  \item Larger $n$ mildly slows convergence: the chain navigates a larger
    neighbourhood ($k(N-k)$ neighbours per state), reducing the probability of
    beneficial swaps.
\end{itemize}

The observed mixing time of $O(10^3)$ steps, even when $|\mathcal{X}|=
\binom{N}{k}$ reaches $\binom{462}{10}\approx 10^{20}$ for $k=10$, is consistent
with the high vertex degree $k(N-k)$ maintaining a large spectral gap.

\subsection{Experiment D: Cost of the Cardinality Constraint}
\label{sec:expD}

Since this experiment studies how the cardinality constraint cost shrinks with
$k$ and does not require a brute-force benchmark, greedy and SA are run on the
full $n=N=462$ universe.  Figure~\ref{fig:gap_vs_k} shows annualised return
(left) and gap relative to the Markowitz upper bound (right) as a function of
portfolio size~$k$.  Table~\ref{tab:gap_vs_k} gives selected values.

\begin{table}[ht]
\centering
\caption{Greedy and SA returns and gaps vs.\ Markowitz UB ($120.24\%$) as $k$
  varies ($n=462$, $\sigma_0=25.64\%$).}
\label{tab:gap_vs_k}
\begin{tabular}{rcccc}
\toprule
$k$ & Greedy & SA & Greedy gap & SA gap \\
\midrule
  1 & $23.30\%$ & $23.30\%$ & $80.62\%$ & $80.62\%$ \\
  2 & $39.66\%$ & $39.66\%$ & $67.01\%$ & $67.01\%$ \\
  3 & $43.48\%$ & $42.67\%$ & $63.84\%$ & $64.51\%$ \\
  5 & $49.67\%$ & $49.67\%$ & $58.69\%$ & $58.69\%$ \\
  7 & $53.85\%$ & $53.96\%$ & $55.22\%$ & $55.12\%$ \\
 10 & $58.92\%$ & $58.47\%$ & $51.00\%$ & $51.37\%$ \\
 15 & $65.25\%$ & $64.07\%$ & $45.73\%$ & $46.71\%$ \\
 20 & $70.08\%$ & $68.37\%$ & $41.72\%$ & $43.14\%$ \\
 30 & $76.76\%$ & $75.93\%$ & $36.16\%$ & $36.85\%$ \\
 50 & $86.90\%$ & $84.30\%$ & $27.73\%$ & $29.89\%$ \\
100 & $101.39\%$ & $97.80\%$ & $15.68\%$ & $18.66\%$ \\
200 & $114.53\%$ & $111.60\%$ &  $4.75\%$ &  $7.19\%$ \\
\bottomrule
\end{tabular}
\end{table}

\begin{figure}[ht]
  \centering
  \includegraphics[width=\linewidth]{fig5_gap_vs_k.png}
  \caption{Return and gap vs.\ Markowitz upper bound as functions of portfolio
    size~$k$ ($n=462$).  The dashed line is $R^*_{\text{MV}}=120.24\%$.}
  \label{fig:gap_vs_k}
\end{figure}

Both greedy and SA monotonically reduce the gap as $k$ grows.  The gap
starts at $80.62\%$ for a single-stock portfolio ($k=1$) and narrows to
$4.75\%$ (greedy) and $7.19\%$ (SA) at $k=200$.  Convergence toward the
Markowitz upper bound is slow relative to the $N=462$ universe: at $k=50$,
the gap still exceeds $27\%$.  Notably, greedy outperforms SA at several
cardinality levels (e.g.\ $k=10$, $15$, $20$), suggesting that 2,000 SA
steps are insufficient to explore the vastly larger neighbourhood
$k(N-k)\approx 4{,}500$ of $\binom{462}{k}$ compared with smaller universes.

\subsection{Experiment E: Constraint Tightness and Problem Hardness}
\label{sec:expE}

We fix $n=35$, $k=5$ ($\binom{35}{5}=324{,}632$ subsets) and sweep $\sigma_0$
from the 5th to the 95th percentile of the empirical distribution of equal-weight
portfolio variances across all subsets.  Figure~\ref{fig:tightness} shows
feasibility fraction, optimal return, and greedy gap as functions of $\sigma_0$.

\begin{figure}[ht]
  \centering
  \includegraphics[width=\linewidth]{fig6_constraint_tightness.png}
  \caption{Feasibility (left), optimal return (middle), and greedy gap (right)
    as functions of $\sigma_0$ ($n=35$, $k=5$).}
  \label{fig:tightness}
\end{figure}

The minimum achievable variance $1/A_S$ for a given subset $S$ can be
substantially lower than its equal-weight variance, since unconstrained weights
allow variance-reducing short positions.  Consequently, a non-trivial fraction of
subsets remain feasible even for tight values of $\sigma_0$.  Empirical hardness,
measured by the greedy gap, peaks at moderate constraint tightness --- where only
a fraction of subsets are feasible and the feasible ones are nearly tied in
return --- and declines at both extremes.

\subsection{Experiment F: Random Sampling Discovery}
\label{sec:expF}

We pre-compute $v(S)$ for all $\binom{35}{5}=324{,}632$ subsets via the
closed-form and simulate random search over 50 trials.
Figure~\ref{fig:discovery} (left) plots the fraction of the optimal return
achieved as a function of the fraction of the search space sampled.
Figure~\ref{fig:discovery} (right) shows the full distribution of $v(S)$ across
feasible subsets.

\begin{figure}[ht]
  \centering
  \includegraphics[width=\linewidth]{fig7_random_sampling.png}
  \caption{Random sampling discovery curve (left) and distribution of all
    feasible portfolio returns (right) for $n=35$, $k=5$.}
  \label{fig:discovery}
\end{figure}

Of the $324{,}632$ subsets, $323{,}027$ ($99.5\%$) are feasible under the chosen
risk budget.  The optimal return across all feasible subsets is $43.03\%$.
Random search reaches $90\%$ of this optimal by sampling only $0.50\%$ of the
search space, and $99\%$ of optimal by sampling $34.17\%$.  The density of
near-optimal portfolios in the search space explains the fast practical
convergence of SA despite $|\mathcal{X}|=\binom{N}{k}$ growing super-exponentially
with $N$.

\subsection{Experiment G: LASSO Relaxation}
\label{sec:expG}

We solve problem~\eqref{eq:lasso} on the same $n=30$ BF\_TICKERS sub-universe
as Experiment~B, sweeping $\lambda$ over 100 log-spaced values from $10^{-3}$
to $1$.  All 30 large-cap stocks in BF\_TICKERS have individual volatility above
the risk limit $\sigma_0=25.64\%$; only four names (V, JNJ, PFE, MRK) can form a
feasible single-stock portfolio.

\paragraph{Regularisation path.}
Figure~\ref{fig:lasso_path} traces the LASSO regularisation path.
\begin{figure}[H]
  \centering
  \includegraphics[width=\linewidth]{fig8_lasso_path.png}
  \caption{LASSO regularisation path ($n=30$, $\sigma_0=25.64\%$).
    \textit{Left}: annualised return vs.\ $\lambda$.
    \textit{Centre}: gross exposure $\|\mathbf{w}\|_1$ vs.\ $\lambda$;
      the red dashed line marks $\|\mathbf{w}\|_1=1$ (long-only).
    \textit{Right}: effective cardinality $k_{\text{eff}}$ (weights $>1\%$ in
      magnitude) vs.\ $\lambda$; the dashed line marks the long-only floor $k=9$.}
  \label{fig:lasso_path}
\end{figure}

At small $\lambda$ the LASSO solution approximates the dense Markowitz portfolio
($k_{\text{eff}}=26$, $\|\mathbf{w}\|_1=4.02$, return $47.79\%$).  As $\lambda$
increases, the penalty progressively shrinks short positions: gross exposure falls
monotonically toward $1$, and $k_{\text{eff}}$ decreases from 26 to~9.  Once
$\|\mathbf{w}\|_1=1$ is reached (all weights non-negative), the path stalls:
$k_{\text{eff}}$ remains at~9 and the return stabilises at $41.31\%$
for all $\lambda \gtrsim 0.06$.

This is a direct empirical manifestation of the \emph{long-only floor} derived in
\S\ref{sec:lasso_theory}: the $\ell_1$ penalty can reduce gross exposure but cannot
reduce cardinality below the minimum long-only feasible count for this universe.
The floor $k_{\min}^{\ell_1}=9$ is a structural property of the data and the risk
constraint, not of the algorithm.

\paragraph{Comparison with cardinality-constrained methods.}
Figure~\ref{fig:lasso_comparison} compares LASSO (at the $k_{\text{eff}}$ values
it achieves) with SA and greedy on the same universe.
Table~\ref{tab:lasso_comparison} quantifies the \emph{price of convexification}:
the return gap between the exact L0 methods and the L1 surrogate at matched
effective cardinality.

\begin{figure}[H]
  \centering
  \includegraphics[width=0.9\linewidth]{fig9_lasso_comparison.png}
  \caption{Method comparison ($n=30$, $\sigma_0=25.64\%$).
    LASSO is plotted only at the $k_{\text{eff}}$ values it achieves ($k\geq 9$);
    the vertical dotted line marks its long-only floor at $k=9$.
    Brute-force stars give the exact optimum for $k\leq 5$.}
  \label{fig:lasso_comparison}
\end{figure}

\begin{table}[ht]
\centering
\caption{Price of convexification: LASSO return vs.\ SA and greedy at matched
  effective cardinality ($n=30$, $\sigma_0=25.64\%$).
  Negative gaps indicate LASSO underperforms the L0 method.}
\label{tab:lasso_comparison}
\begin{tabular}{rcccrr}
\toprule
$k$ & SA & Greedy & LASSO & LASSO$-$SA & LASSO$-$Greedy \\
\midrule
 9 & $44.62\%$ & $42.70\%$ & $41.31\%$ & $-7.40\%$ & $-3.25\%$ \\
10 & $45.18\%$ & $43.68\%$ & $41.45\%$ & $-8.25\%$ & $-5.10\%$ \\
12 & $45.82\%$ & $45.60\%$ & $41.73\%$ & $-8.91\%$ & $-8.48\%$ \\
15 & $46.71\%$ & $46.46\%$ & $43.67\%$ & $-6.51\%$ & $-6.02\%$ \\
17 & $47.19\%$ & $46.96\%$ & $45.22\%$ & $-4.16\%$ & $-3.70\%$ \\
19 & $47.51\%$ & $47.40\%$ & $46.41\%$ & $-2.31\%$ & $-2.09\%$ \\
\bottomrule
\end{tabular}
\end{table}

Two observations stand out.  First, LASSO consistently underperforms both SA
and greedy at every matched cardinality level.  The gap is largest at $k=12$
($-8.91\%$ vs.\ SA) and narrows as $k$ grows toward the Markowitz UB ($47.79\%$),
because the L1 penalty has less work to do when the unconstrained solution is already
close to long-only.  Second, the LASSO floor at $k=9$ ($41.31\%$) lies below the
brute-force optimum at $k=5$ ($41.79\%$), yet uses $80\%$ more stocks.
This illustrates the key cost of the L1 relaxation: it cannot replicate the
short-selling strategy that allows exact methods to achieve high returns with few
stocks.  The exact $k=5$ solution (NVDA, JPM, WFC, JNJ, LLY) uses a mix of
long and short positions ($\|\mathbf{w}\|_1>1$) that is penalised out of the
LASSO path.

% ─────────────────────────────────────────────────────────────────────────────
\section{Discussion}
\label{sec:discussion}
% ─────────────────────────────────────────────────────────────────────────────

\paragraph{Tractability of the inner problem.}
Because short selling is permitted, the only constraints on $\mathbf{w}_S$ are
the budget equation $\mathbf{1}^\top\mathbf{w}_S=1$ and the variance bound.
The inner problem~\eqref{eq:inner} is therefore an equality-constrained strictly
convex quadratic program --- exactly the classical Markowitz mean-variance
problem --- and its solution is given in closed form by
equations~\eqref{eq:mvreturn}--\eqref{eq:mvweights}.  Each evaluation of $v(S)$
reduces to a single $k\times k$ matrix inversion and a few dot products,
incurring no iterative solver overhead.  The sole computational bottleneck of
problem~\eqref{eq:card} therefore lies in the combinatorial outer search over
$\binom{n}{k}$ subsets, not in the per-subset optimisation.

\paragraph{Short selling and feasibility.}
The minimum achievable variance for a $k$-subset $S$ is $\sigma^2_{\min,S}=1/A_S$,
attained by the global minimum-variance portfolio of $S$.  Because weights are
unconstrained, short positions allow the variance to be pushed well below what
any convex combination of the assets could achieve, so the feasibility threshold
$\sigma_0^2 \geq 1/A_S$ is met for a wide range of subsets.  Concretely,
Experiment~F establishes that $99.5\%$ of all $\binom{35}{5}=324{,}632$
five-asset subsets are feasible under the calibrated risk budget $\sigma_0=25.64\%$.

\paragraph{NP-hardness of the outer search.}
The tractability of the inner problem does not extend to the cardinality-constrained
problem~\eqref{eq:card}.  The combinatorial state space $\mathcal{X}=\binom{[N]}{k}$
has cardinality $\binom{N}{k}$, which grows super-exponentially with $N$; no
polynomial-time algorithm is known for the outer optimisation.  Greedy and SA
therefore remain heuristics, and their quality must be assessed empirically
against the brute-force optimum on tractable instances.

\paragraph{Risk budget calibration.}
Equation~\eqref{eq:riskbudget} anchors $\sigma_0^2$ to individual stock
variances, not portfolio variance.  This is essential: calibrating on the
diversified $N$-stock portfolio variance ($\approx 3.26\%$ annualised) would
produce a limit far below $\sigma^2_{\min,S}$ for small $k$, rendering every
inner problem infeasible.

\paragraph{L1 relaxation and the long-only floor.}
Experiment~G reveals that the L1 penalty in the LASSO surrogate~\eqref{eq:lasso}
operates primarily by eliminating short positions (reducing $\|\mathbf{w}\|_1$
from 4.02 to 1), not by directly selecting fewer stocks.  Once the solution is
long-only, further sparsification via $\ell_1$ is structurally impossible (the
floor at $k=9$), and the price of convexification is substantial: LASSO at $k=9$
yields $41.31\%$, below SA at $k=5$ ($41.79\%$) which uses short positions to
concentrate return in fewer stocks.  These results suggest that the L1 surrogate
is most appropriate when short selling is restricted by mandate, or when gross
exposure is the primary concern rather than stock-count reduction.

\paragraph{Estimation risk.}
With $T/N\approx 6$, the full $462\times 462$ sample covariance matrix
is near-rank-deficient, so the Markowitz upper bound of $120.24\%$ reflects
in-sample optimisation under noisy estimates rather than a meaningful out-of-sample
target (see \S\ref{sec:rmt}).  The cardinality-constrained results for $k\leq 30$
are unaffected, as the corresponding $k\times k$ sub-matrices satisfy $T/k>90$.
Shrinkage estimators (e.g.\ Ledoit--Wolf \cite{LedoitWolf2004}) or factor models
would improve the conditioning of $\hat{\boldsymbol{\Sigma}}$ and yield more
reliable estimates of the Markowitz upper bound; this is a natural direction for
future work.

\paragraph{Interpreting the Markowitz upper bound.}
The Markowitz upper bound of $R^*_{\text{MV}}=120.24\%$ should not be
read as a realistic performance target.
It is an in-sample artefact produced by inverting a $462\times462$
sample covariance matrix whose smallest eigenvalues are dominated by
estimation noise (see \S\ref{sec:rmt}): the optimizer exploits these
noise directions to construct extreme long--short positions
(231 long, 224 short) that look profitable in-sample but would not
generalise out-of-sample.
Its role in our analysis is purely as a \emph{ceiling} that anchors the
gap curves in Experiments~D and~G since it tells us how much return is
theoretically available to an unconstrained optimizer on the same data,
not how much an investor could realistically earn.
The comparisons that carry practical weight are those \emph{among} the
cardinality-constrained methods: SA versus greedy on the same
sub-universe (Experiment~B), the price of convexification between L0
and L1 methods at matched cardinality (Experiment~G), and the
convergence behaviour of SA across different values of $k$ (Experiment~C).
These comparisons are unaffected by the full-$N$ conditioning issue
because they operate on $k\times k$ sub-matrices.
A shrinkage estimator applied
to the full covariance matrix would compress the upper bound toward a
more credible figure, but would leave the relative rankings among
cardinality-constrained methods unchanged, since those rankings depend
on the sub-matrix structure rather than the full inverse.

% ─────────────────────────────────────────────────────────────────────────────
\section{Conclusion}
\label{sec:conclusion}
% ─────────────────────────────────────────────────────────────────────────────

We studied cardinality-constrained Markowitz portfolio optimisation on $N=462$
S\&P~500 stocks ($T=2{,}854$ trading days), with short selling allowed.
The key findings are:
\begin{enumerate}[leftmargin=1.5em]
  \item Without a cardinality constraint the problem reduces to classical
    Markowitz mean-variance optimisation and has a closed-form solution via
    the two-fund separation theorem, achieving $R^*_{\text{MV}}=120.24\%$
    (in-sample) at $\sigma_0=25.64\%$ with 231 long and 224 short positions.
  \item Adding $|\operatorname{supp}(\mathbf{w})|=k$ makes the outer search
    super-exponential; the closed-form inner solution is used to evaluate
    each of the $\binom{N}{k}$ candidate subsets efficiently.
  \item Brute force becomes infeasible around $n\approx 40$--45 for $k=5$;
    on the curated 30-stock sub-universe, SA matches the brute-force optimal
    exactly ($41.79\%$, $0.00\%$ gap) while greedy obtains $40.49\%$
    ($3.11\%$ gap) at $k=5$.
  \item SA, viewed as a Markov chain on $\binom{[N]}{k}$, converges within
    $\sim\!500$--$1{,}000$ steps even for $|\mathcal{X}|\approx 10^{20}$
    ($N=462$, $k=10$), consistent with a large spectral gap due to high vertex
    degree $k(N-k)$.
  \item The cardinality constraint gap decreases monotonically with $k$, from
    $80.62\%$ at $k=1$ to $4.75\%$ at $k=200$ (greedy); convergence is
    noticeably slower for large universes.
  \item Empirical hardness (greedy gap) peaks at moderate constraint tightness.
  \item $99.5\%$ of all 5-subsets of 35 stocks are feasible; random search
    achieves $99\%$ of the optimal return by sampling $34.17\%$ of the search
    space, confirming a dense return landscape near the optimum.
  \item The L1 (LASSO) relaxation of the cardinality constraint is solvable as
    an SOCP but reveals a structural limitation: it reduces cardinality by
    eliminating short positions (gross exposure $\|\mathbf{w}\|_1$ falls from
    4.02 to~1), not by directly selecting fewer stocks.  Once the portfolio is
    long-only, further sparsification is impossible via $\ell_1$ alone.  The
    resulting \emph{long-only floor} at $k=9$ yields $41.31\%$, below the
    exact $k=5$ optimum ($41.79\%$) that exploits short selling.
    The price of convexification at matched cardinality ranges from $-2.3\%$
    to $-8.9\%$ vs.\ SA.
\end{enumerate}

% ─────────────────────────────────────────────────────────────────────────────
\begin{thebibliography}{9}
% ─────────────────────────────────────────────────────────────────────────────

\bibitem{Markowitz1952}
Markowitz, H.\ (1952).
Portfolio selection.
\textit{Journal of Finance}, 7(1), 77--91.

\bibitem{GareyJohnson1979}
Garey, M.~R.\ and Johnson, D.~S.\ (1979).
\textit{Computers and Intractability: A Guide to the Theory of
NP-Completeness}.
W.~H.\ Freeman, New York.

\bibitem{MarchenkoPastur1967}
Mar\v{c}enko, V.~A.\ and Pastur, L.~A.\ (1967).
Distribution of eigenvalues for some sets of random matrices.
\textit{Mathematics of the USSR-Sbornik}, 1(4), 457--483.

\bibitem{Tibshirani1996}
Tibshirani, R.\ (1996).
Regression shrinkage and selection via the lasso.
\textit{Journal of the Royal Statistical Society: Series B}, 58(1), 267--288.

\bibitem{LedoitWolf2004}
Ledoit, O.\ and Wolf, M.\ (2004).
A well-conditioned estimator for large-dimensional covariance matrices.
\textit{Journal of Multivariate Analysis}, 88(2), 365--411.

\bibitem{Bienstock1996}
Bienstock, D.\ (1996).
Computational study of a family of mixed-integer quadratic programming
problems.
\textit{Mathematical Programming}, 74(2), 121--140.

\bibitem{Godsil1993}
Godsil, C.\ (1993).
\textit{Algebraic Combinatorics}.
Chapman \& Hall, New York.

\bibitem{LevinPeres2017}
Levin, D.~A.\ and Peres, Y.\ (2017).
\textit{Markov Chains and Mixing Times} (2nd ed.).
American Mathematical Society, Providence.

\bibitem{LawlerSokal1988}
Lawler, G.~F.\ and Sokal, A.~D.\ (1988).
Bounds on the $L^2$ spectrum for Markov chains and Markov processes:
a generalization of Cheeger's inequality.
\textit{Transactions of the American Mathematical Society},
309(2), 557--580.

\bibitem{KellererMansiniSperanza2000}
Kellerer, H., Mansini, R., and Speranza, M.~G.\ (2000).
Selecting portfolios with fixed costs and minimum transaction lots.
\textit{Annals of Operations Research}, 99(1--4), 287--304.

\end{thebibliography}

\end{document}